\documentclass[10pt,a4paper]{amsart}
\usepackage{amsmath,amssymb,amsthm,graphicx}
% Retained preamble for the Stacks Project chapter snapshot in this directory.
% The chapter's own preamble.tex is a copy of the shared Stacks Project preamble; the block
% below is sliced MECHANICALLY, line for line, from it. Nothing was added, removed or
% redefined; the wrapper class, the xr-hyper/hyperref wiring and the externaldocument list
% are omitted because the assembler supplies the document class.
%
% lines taken from preamble.tex: 3, 5, 9, 11, 12, 156, 157, 158, ... (51 lines total)
%
% package coverage restored from the chapter's own preamble (2026-10-13): xy, xy-2cell

\usepackage[all]{xy}
\xyoption{2cell}
\UseAllTwocells

\newenvironment{reference}{\comment}{\endcomment}
\newenvironment{slogan}{\comment}{\endcomment}
\newenvironment{history}{\comment}{\endcomment}
\theoremstyle{plain}
\newtheorem{theorem}[subsection]{Theorem}
\newtheorem{proposition}[subsection]{Proposition}
\newtheorem{lemma}[subsection]{Lemma}
\theoremstyle{definition}
\newtheorem{definition}[subsection]{Definition}
\newtheorem{example}[subsection]{Example}
\newtheorem{exercise}[subsection]{Exercise}
\newtheorem{situation}[subsection]{Situation}
\theoremstyle{remark}
\newtheorem{remark}[subsection]{Remark}
\newtheorem{remarks}[subsection]{Remarks}
\numberwithin{equation}{subsection}
\def\lim{\mathop{\mathrm{lim}}\nolimits}
\def\colim{\mathop{\mathrm{colim}}\nolimits}
\def\Spec{\mathop{\mathrm{Spec}}}
\def\Hom{\mathop{\mathrm{Hom}}\nolimits}
\def\Ext{\mathop{\mathrm{Ext}}\nolimits}
\def\SheafHom{\mathop{\mathcal{H}\!\mathit{om}}\nolimits}
\def\SheafExt{\mathop{\mathcal{E}\!\mathit{xt}}\nolimits}
\def\Sch{\mathit{Sch}}
\def\Mor{\mathop{\mathrm{Mor}}\nolimits}
\def\Ob{\mathop{\mathrm{Ob}}\nolimits}
\def\Sh{\mathop{\mathit{Sh}}\nolimits}
\def\NL{\mathop{N\!L}\nolimits}
\def\CH{\mathop{\mathrm{CH}}\nolimits}
\def\proetale{{pro\text{-}\acute{e}tale}}
\def\etale{{\acute{e}tale}}
\def\QCoh{\mathit{QCoh}}
\def\Ker{\mathop{\mathrm{Ker}}}
\def\Im{\mathop{\mathrm{Im}}}
\def\Coker{\mathop{\mathrm{Coker}}}
\def\Coim{\mathop{\mathrm{Coim}}}
\DeclareMathSymbol{\boxtimes}{\mathbin}{AMSa}{"02}
\def\QCohstack{\mathcal{QC}\!\mathit{oh}}
\def\Cohstack{\mathcal{C}\!\mathit{oh}}
\def\Spacesstack{\mathcal{S}\!\mathit{paces}}
\def\Quotfunctor{\mathrm{Quot}}
\def\Hilbfunctor{\mathrm{Hilb}}
\def\Curvesstack{\mathcal{C}\!\mathit{urves}}
\def\Polarizedstack{\mathcal{P}\!\mathit{olarized}}
\def\Complexesstack{\mathcal{C}\!\mathit{omplexes}}
\def\Pic{\mathop{\mathrm{Pic}}\nolimits}
\def\Picardstack{\mathcal{P}\!\mathit{ic}}
\def\Picardfunctor{\mathrm{Pic}}
\def\Deformationcategory{\mathcal{D}\!\mathit{ef}}
\usepackage{cleveref}
\crefname{originallemma}{original lemma}{original lemmas}
\Crefname{originallemma}{Original lemma}{Original lemmas}
\makeatletter
\newlabel{algebra-lemma-descend-properties-modules}{{03C4}{source}{}{sourcepointer-1}{}}
\newlabel{algebra-lemma-descend-properties-modules@cref}{{[originallemma][0][]03C4}{[1][0][]source}{}{}{}}
\makeatother
\crefname{originallemma}{original lemma}{original lemmas}
\Crefname{originallemma}{Original lemma}{Original lemmas}
\makeatletter
\newlabel{lemma-quasi-coherent-strongly-cartesian}{{04U3}{source}{}{sourcepointer-2}{}}
\newlabel{lemma-quasi-coherent-strongly-cartesian@cref}{{[originallemma][0][]04U3}{[1][0][]source}{}{}{}}
\makeatother
\crefname{originallemma}{original lemma}{original lemmas}
\Crefname{originallemma}{Original lemma}{Original lemmas}
\makeatletter
\newlabel{modules-lemma-pullback-finite-type}{{01B6}{source}{}{sourcepointer-3}{}}
\newlabel{modules-lemma-pullback-finite-type@cref}{{[originallemma][0][]01B6}{[1][0][]source}{}{}{}}
\makeatother
\crefname{originallemma}{original lemma}{original lemmas}
\Crefname{originallemma}{Original lemma}{Original lemmas}
\makeatletter
\newlabel{stacks-definition-stack}{{026F}{source}{}{sourcepointer-4}{}}
\newlabel{stacks-definition-stack@cref}{{[originallemma][0][]026F}{[1][0][]source}{}{}{}}
\makeatother
\crefname{originallemma}{original lemma}{original lemmas}
\Crefname{originallemma}{Original lemma}{Original lemmas}
\makeatletter
\newlabel{stacks-lemma-substack}{{04TU}{source}{}{sourcepointer-5}{}}
\newlabel{stacks-lemma-substack@cref}{{[originallemma][0][]04TU}{[1][0][]source}{}{}{}}
\makeatother
\title{Finite-type quasi-coherent modules form a stack}
\author{The Stacks Project Authors}
\date{stacks-project:examples-stacks.tex:local-snapshot-20261012; GNU FDL 1.2 or later}
\begin{document}
\maketitle
\noindent\textit{Editorial note.} Statement span is lines 416-421 of sources/stacks-examples-stacks/pinned-original-source.tex; proof span is lines 423-447, the complete author proof (23 body lines).\par\medskip
\noindent\textit{Editorial note.} Navigation stubs: 5 cited-but-not-included labels are supplied as non-load-bearing pointers with their pinned tags where the snapshot tag index has them.\par\medskip
\noindent\textit{Editorial note.} extra\_packages is empty on purpose: the retained preamble already loads xy with options and a second declaration clashes in the assembler wrapper.\par\medskip
\noindent\textit{Editorial note.} Text is retained verbatim, including the author's own typos and placeholders.\par\medskip
\noindent\textit{Editorial note.} Queue index 10 of queue\_examples-stacks.json (pb=1218, stmt 416-421, proof 423-447, halfA).\par\medskip
\noindent\textit{Editorial note.} stacks-examples-stacks wave-1 (exs1): fresh adjudication accept (2026-10-14); pool per Chief 09:4x authorization (idle-line throughput over kun's backup note, SMC precedent); examples-stacks.tex 27 labels/20 with proof judged 9 accept / 11 reject; claim\_key v2 stacks:examples-stacks.tex:<label>.\par\medskip
\noindent\textit{Allowed standard prerequisites (editorial).}\begin{itemize}
\item Explicit prerequisite boundary. The cited inputs are prerequisites; the argument in the retained span is not replaced by any cited result.
\end{itemize}\medskip
\ifdefined\hypertarget\hypertarget{sourcepointer-1}{}\fi
\noindent\textit{Source reference.} Original lemma 03C4: sources/stacks-examples-stacks/pinned-original-source.tex, source lines 447--447. Cited by the retained statement or proof but not part of the collected spans; supplied only as a navigation pointer. The argument does not depend on its content.\par\medskip
\ifdefined\hypertarget\hypertarget{sourcepointer-2}{}\fi
\noindent\textit{Source reference.} Original lemma 04U3: sources/stacks-examples-stacks/pinned-original-source.tex, source lines 333--337. Cited by the retained statement or proof but not part of the collected spans; supplied only as a navigation pointer. The argument does not depend on its content.\par\medskip
\ifdefined\hypertarget\hypertarget{sourcepointer-3}{}\fi
\noindent\textit{Source reference.} Original lemma 01B6: sources/stacks-examples-stacks/pinned-original-source.tex, source lines 447--447. Cited by the retained statement or proof but not part of the collected spans; supplied only as a navigation pointer. The argument does not depend on its content.\par\medskip
\ifdefined\hypertarget\hypertarget{sourcepointer-4}{}\fi
\noindent\textit{Source reference.} Original lemma 026F: sources/stacks-examples-stacks/pinned-original-source.tex, source lines 447--447. Cited by the retained statement or proof but not part of the collected spans; supplied only as a navigation pointer. The argument does not depend on its content.\par\medskip
\ifdefined\hypertarget\hypertarget{sourcepointer-5}{}\fi
\noindent\textit{Source reference.} Original lemma 04TU: sources/stacks-examples-stacks/pinned-original-source.tex, source lines 447--447. Cited by the retained statement or proof but not part of the collected spans; supplied only as a navigation pointer. The argument does not depend on its content.\par\medskip
\begin{lemma}
\label{lemma-stack-of-finite-type-quasi-coherent-sheaves}
The functor $p_{fg} : \QCohstack_{fg} \to (\Sch/S)_{fppf}$
satisfies conditions (1), (2) and (3) of
Stacks, Definition \ref{stacks-definition-stack}.
\end{lemma}

\begin{proof}
We will verify assumptions (1), (2), (3) of
Stacks, Lemma \ref{stacks-lemma-substack}
to prove this. By
Lemma \ref{lemma-quasi-coherent-strongly-cartesian}
a morphism $(Y, \mathcal{G}) \to (X, \mathcal{F})$ is
strongly cartesian if and only if it induces an isomorphism
$f^*\mathcal{F} \to \mathcal{G}$. By
Modules, Lemma \ref{modules-lemma-pullback-finite-type}
the pullback of a finite type $\mathcal{O}_X$-module is of finite
type. Hence assumption (1) of
Stacks, Lemma \ref{stacks-lemma-substack}
holds. Assumption (2) holds trivially.
Finally, to prove assumption (3) we have to show:
If $\mathcal{F}$ is a quasi-coherent $\mathcal{O}_X$-module
and $\{f_i : X_i \to X\}$ is an fppf covering such that each
$f_i^*\mathcal{F}$ is of finite type, then $\mathcal{F}$ is of
finite type. Considering the restriction of $\mathcal{F}$ to
an affine open of $X$ this reduces to the following algebra statement:
Suppose that $R \to S$ is a finitely presented, faithfully flat ring map
and $M$ an $R$-module. If $M \otimes_R S$ is a finitely generated
$S$-module, then $M$ is a finitely generated $R$-module.
A stronger form of the algebra fact can be found in
Algebra, Lemma \ref{algebra-lemma-descend-properties-modules}.
\end{proof}

\end{document}
